\documentclass[a4paper,11pt,oneside]{memoir}

% Castellano
\usepackage[spanish,es-tabla]{babel}
\selectlanguage{spanish}
\usepackage[utf8]{inputenc}
\usepackage{placeins}

\RequirePackage{booktabs}
\RequirePackage[table]{xcolor}
\RequirePackage{xtab}
\RequirePackage{multirow}

% Links
\usepackage[colorlinks]{hyperref}
\hypersetup{
	allcolors = {red}
}

% Ecuaciones
\usepackage{amsmath}

% Rutas de fichero / paquete
\newcommand{\ruta}[1]{{\sffamily #1}}

% Párrafos
\nonzeroparskip


% Imagenes
\usepackage{graphicx}
\newcommand{\imagen}[2]{
	\begin{figure}[!h]
		\centering
		\includegraphics[width=0.9\textwidth]{#1}
		\caption{#2}\label{fig:#1}
	\end{figure}
	\FloatBarrier
}

\newcommand{\imagenflotante}[2]{
	\begin{figure}%[!h]
		\centering
		\includegraphics[width=0.9\textwidth]{#1}
		\caption{#2}\label{fig:#1}
	\end{figure}
}



% El comando \figura nos permite insertar figuras comodamente, y utilizando
% siempre el mismo formato. Los parametros son:
% 1 -> Porcentaje del ancho de página que ocupará la figura (de 0 a 1)
% 2 --> Fichero de la imagen
% 3 --> Texto a pie de imagen
% 4 --> Etiqueta (label) para referencias
% 5 --> Opciones que queramos pasarle al \includegraphics
% 6 --> Opciones de posicionamiento a pasarle a \begin{figure}
\newcommand{\figuraConPosicion}[6]{%
  \setlength{\anchoFloat}{#1\textwidth}%
  \addtolength{\anchoFloat}{-4\fboxsep}%
  \setlength{\anchoFigura}{\anchoFloat}%
  \begin{figure}[#6]
    \begin{center}%
      \Ovalbox{%
        \begin{minipage}{\anchoFloat}%
          \begin{center}%
            \includegraphics[width=\anchoFigura,#5]{#2}%
            \caption{#3}%
            \label{#4}%
          \end{center}%
        \end{minipage}
      }%
    \end{center}%
  \end{figure}%
}

%
% Comando para incluir imágenes en formato apaisado (sin marco).
\newcommand{\figuraApaisadaSinMarco}[5]{%
  \begin{figure}%
    \begin{center}%
    \includegraphics[angle=90,height=#1\textheight,#5]{#2}%
    \caption{#3}%
    \label{#4}%
    \end{center}%
  \end{figure}%
}
% Para las tablas
\newcommand{\otoprule}{\midrule [\heavyrulewidth]}
%
% Nuevo comando para tablas pequeñas (menos de una página).
\newcommand{\tablaSmall}[5]{%
 \begin{table}
  \begin{center}
   \rowcolors {2}{gray!35}{}
   \begin{tabular}{#2}
    \toprule
    #4
    \otoprule
    #5
    \bottomrule
   \end{tabular}
   \caption{#1}
   \label{tabla:#3}
  \end{center}
 \end{table}
}

%
% Nuevo comando para tablas pequeñas (menos de una página).
\newcommand{\tablaSmallSinColores}[5]{%
 \begin{table}[H]
  \begin{center}
   \begin{tabular}{#2}
    \toprule
    #4
    \otoprule
    #5
    \bottomrule
   \end{tabular}
   \caption{#1}
   \label{tabla:#3}
  \end{center}
 \end{table}
}

\newcommand{\tablaApaisadaSmall}[5]{%
\begin{landscape}
  \begin{table}
   \begin{center}
    \rowcolors {2}{gray!35}{}
    \begin{tabular}{#2}
     \toprule
     #4
     \otoprule
     #5
     \bottomrule
    \end{tabular}
    \caption{#1}
    \label{tabla:#3}
   \end{center}
  \end{table}
\end{landscape}
}

%
% Nuevo comando para tablas grandes con cabecera y filas alternas coloreadas en gris.
\newcommand{\tabla}[6]{%
  \begin{center}
    \tablefirsthead{
      \toprule
      #5
      \otoprule
    }
    \tablehead{
      \multicolumn{#3}{l}{\small\sl continúa desde la página anterior}\\
      \toprule
      #5
      \otoprule
    }
    \tabletail{
      \hline
      \multicolumn{#3}{r}{\small\sl continúa en la página siguiente}\\
    }
    \tablelasttail{
      \hline
    }
    \bottomcaption{#1}
    \rowcolors {2}{gray!35}{}
    \begin{xtabular}{#2}
      #6
      \bottomrule
    \end{xtabular}
    \label{tabla:#4}
  \end{center}
}

%
% Nuevo comando para tablas grandes con cabecera.
\newcommand{\tablaSinColores}[6]{%
  \begin{center}
    \tablefirsthead{
      \toprule
      #5
      \otoprule
    }
    \tablehead{
      \multicolumn{#3}{l}{\small\sl continúa desde la página anterior}\\
      \toprule
      #5
      \otoprule
    }
    \tabletail{
      \hline
      \multicolumn{#3}{r}{\small\sl continúa en la página siguiente}\\
    }
    \tablelasttail{
      \hline
    }
    \bottomcaption{#1}
    \begin{xtabular}{#2}
      #6
      \bottomrule
    \end{xtabular}
    \label{tabla:#4}
  \end{center}
}

%
% Nuevo comando para tablas grandes sin cabecera.
\newcommand{\tablaSinCabecera}[5]{%
  \begin{center}
    \tablefirsthead{
      \toprule
    }
    \tablehead{
      \multicolumn{#3}{l}{\small\sl continúa desde la página anterior}\\
      \hline
    }
    \tabletail{
      \hline
      \multicolumn{#3}{r}{\small\sl continúa en la página siguiente}\\
    }
    \tablelasttail{
      \hline
    }
    \bottomcaption{#1}
  \begin{xtabular}{#2}
    #5
   \bottomrule
  \end{xtabular}
  \label{tabla:#4}
  \end{center}
}



\definecolor{cgoLight}{HTML}{EEEEEE}
\definecolor{cgoExtralight}{HTML}{FFFFFF}

%
% Nuevo comando para tablas grandes sin cabecera.
\newcommand{\tablaSinCabeceraConBandas}[5]{%
  \begin{center}
    \tablefirsthead{
      \toprule
    }
    \tablehead{
      \multicolumn{#3}{l}{\small\sl continúa desde la página anterior}\\
      \hline
    }
    \tabletail{
      \hline
      \multicolumn{#3}{r}{\small\sl continúa en la página siguiente}\\
    }
    \tablelasttail{
      \hline
    }
    \bottomcaption{#1}
    \rowcolors[]{1}{cgoExtralight}{cgoLight}

  \begin{xtabular}{#2}
    #5
   \bottomrule
  \end{xtabular}
  \label{tabla:#4}
  \end{center}
}


















\graphicspath{ {./img/} }

% Capítulos
\chapterstyle{bianchi}
\newcommand{\capitulo}[2]{
	\setcounter{chapter}{#1}
	\setcounter{section}{0}
	\chapter*{#2}
	\addcontentsline{toc}{chapter}{#2}
	\markboth{#2}{#2}
}

% Apéndices
\renewcommand{\appendixname}{Apéndice}
\renewcommand*\cftappendixname{\appendixname}

\newcommand{\apendice}[1]{
	%\renewcommand{\thechapter}{A}
	\chapter{#1}
}

\renewcommand*\cftappendixname{\appendixname\ }

% Formato de portada
\makeatletter
\usepackage{xcolor}
\newcommand{\tutor}[1]{\def\@tutor{#1}}
\newcommand{\contutor}[1]{\def\@contutor{#1}}
\newcommand{\course}[1]{\def\@course{#1}}
\definecolor{cpardoBox}{HTML}{E6E6FF}
\def\maketitle{
  \null
  \thispagestyle{empty}
  % Cabecera ----------------
\noindent\includegraphics[width=\textwidth]{cabecera}\vspace{1cm}%
  \vfill
  % Título proyecto y escudo informática ----------------
  \colorbox{cpardoBox}{%
    \begin{minipage}{.8\textwidth}
      \vspace{.5cm}\Large
      \begin{center}
      \textbf{TFG del Grado en Ingeniería Informática}\vspace{.6cm}\\
      \textbf{\LARGE\@title{}}
      \end{center}
      \vspace{.2cm}
    \end{minipage}

  }%
  \hfill\begin{minipage}{.20\textwidth}
    \includegraphics[width=\textwidth]{escudoInfor}
  \end{minipage}
  \vfill
  % Datos de alumno, curso y tutores ------------------
  \begin{center}%
  {%
    \noindent\LARGE
    Presentado por \@author{}\\ 
    en Universidad de Burgos --- \@date{}\\
    Tutor: \@tutor{}\\
    Contutor: \@contutor{}\\
  }%
  \end{center}%
  \null
  \cleardoublepage
  }
\makeatother

\newcommand{\nombre}{Segundo González González} %%% cambio de comando
\newcommand{\dni}{20.219.016-S}
\newcommand{\tutorA}{Carlos López Nozal}
\newcommand{\tutorB}{José Francisco Díez Pastor}
\newcommand{\dptotutorA}{Ingeniería Informática}
\newcommand{\areatutorA}{Lenguajes y Sistemas Informáticos}
\newcommand{\titulo}{Investigación, creación y explotación de modelos para el estudio del COVID-19 en pacientes del Servicio Cántabro de Salud}

% Datos de portada
\title{\titulo}
\author{\nombre}
\tutor{\tutorA}
\contutor{\tutorB}
\date{\today}

\begin{document}

\maketitle


\newpage\null\thispagestyle{empty}\newpage


%%%%%%%%%%%%%%%%%%%%%%%%%%%%%%%%%%%%%%%%%%%%%%%%%%%%%%%%%%%%%%%%%%%%%%%%%%%%%%%%%%%%%%%%
\thispagestyle{empty}


\noindent\includegraphics[width=\textwidth]{cabecera}\vspace{1cm}

\noindent D. \tutorA, profesor del departamento de \dptotutorA, área de \areatutorA.

\noindent Expone:

\noindent Que el alumno D. \nombre, con DNI \dni, ha realizado el Trabajo final de Grado en Ingeniería Informática titulado "\titulo". 

\noindent Y que dicho trabajo ha sido realizado por el alumno bajo la dirección del que suscribe, en virtud de lo cual se autoriza su presentación y defensa.

\begin{center} %\large
En Burgos, {\large \today}
\end{center}

\vfill\vfill\vfill

% Author and supervisor
\begin{minipage}{0.45\textwidth}
\begin{flushleft} %\large
Vº. Bº. del Tutor:\\[2cm]
D. \tutorA
\end{flushleft}
\end{minipage}
\hfill
\begin{minipage}{0.45\textwidth}
\begin{flushleft} %\large
Vº. Bº. del Contutor:\\[2cm]
D. \tutorB
\end{flushleft}
\end{minipage}
\hfill

\vfill

% para casos con solo un tutor comentar lo anterior
% y descomentar lo siguiente
%Vº. Bº. del Tutor:\\[2cm]
%D. nombre tutor


\newpage\null\thispagestyle{empty}\newpage




\frontmatter

% Abstract en castellano
\renewcommand*\abstractname{Resumen}
\begin{abstract}
Para este TFG se cuenta con datos proporcionados por el Instituto de Investigación sanitaria de Valdecilla. El objetivo de este trabajo será la aplicación de técnicas de minería de datos que proporcionen modelos interpretables por el facultativo médico en relación al tratamiento de pacientes con COVID-19.
\end{abstract}

\renewcommand*\abstractname{Descriptores}
\begin{abstract}
COVID-19, machine learning, python \ldots
\end{abstract}

\clearpage

% Abstract en inglés
\renewcommand*\abstractname{Abstract}
\begin{abstract}
For this TFG there are data provided by the Valdecilla Health Research Institute. The objective of this work will be the application of data mining techniques that provide models that can be interpreted by the medical practitioner in relation to the treatment of patients with COVID-19.
\end{abstract}

\renewcommand*\abstractname{Keywords}
\begin{abstract}
COVID-19, machine learning, python \ldots
\end{abstract}

\clearpage

% Indices
\tableofcontents

\clearpage

\listoffigures

\clearpage

\listoftables
\clearpage

\mainmatter
\capitulo{1}{Introducción}

Descripción del contenido del trabajo y del estructura de la memoria y del resto de materiales entregados.

\capitulo{2}{Objetivos del proyecto}

Este apartado explica de forma precisa y concisa cuales son los objetivos que se persiguen con la realización del proyecto. Se puede distinguir entre los objetivos marcados por los requisitos del software a construir y los objetivos de carácter técnico que plantea a la hora de llevar a la práctica el proyecto.

\capitulo{3}{Conceptos teóricos}

\section{Minería de Datos}

La minería de datos es un proceso de exploración y análisis de grandes conjuntos de datos con el objetivo de descubrir patrones, relaciones y conocimientos útiles. Se utiliza en diversas industrias y disciplinas para tomar decisiones basadas en datos y obtener información valiosa. La cantidad de datos que se produce cada día crece de forma exponencial en todos los aspectos de la vida. Esto nos ha llevado a utilizar técnicas como el aprendizaje automático para poder explotar esos datos y sacarles un rendimiento.

Para poder extraer conocimiento de esta ingente cantidad de datos existen diversas metodologías como CRISP-DM, Scrum, Kanban... Nosotros hemos utilizado la primera para crear los modelos de aprendizaje automático.

\section{CRISP-DM}

La metodología CRISP-DM (Cross-Industry Standard Process for Data Mining) es un enfoque utilizado para abordar proyectos de minería de datos. Aunque originalmente se desarrolló para aplicaciones comerciales, se puede adaptar para analizar y abordar los desafíos relacionados con el COVID-19.

Se compone de 6 fases:
\begin{itemize}
	\item \textbf{Entendimiento del Negocio} Comprensión clara de los objetivos y requisitos relacionados con el proyecto.
	\item \textbf{Entendimiento de los Datos} Recopilación y exploración de datos para comprender su calidad, estructura y características, identificando posibles problemas o limitaciones que nos podrían afectar a los análisis posteriores.
    \item \textbf{Preparación de los datos} Realización de tareas de limpieza y transformación de los datos, imputando datos que pudieran faltar, eliminando duplicados, normalizando dichos datos y seleccionando características relevantes.
    \item \textbf{Modelado} y análisis de datos para responder a las preguntas planteadas en la etapa de comprensión del negocio, incluyendo la construcción de modelos predictivos. En esta fase puede ser necesario volver a la fase anterior para volver a preparar los datos.
    \item \textbf{Evaluación} de los modelos desarrollados, analizando los resultados y comparándolos con los objetivos iniciales. Se considera la precisión, la eficacia y la calidad de los modelos.
    \item \textbf{Despliegue} del modelo construido que haya dado mejor resultado para que pueda ser utilizado con otros datos.
\end{itemize}

\imagen{crisp-dm}{Metodología de datos CRISP-DM \cite{crisp-img:2021}}

A continuación, se proporciona un resumen de cómo se ha aplicado la metodología CRISP-DM a este proyecto:

\subsection{Entendimiento del negocio}
El triaje sanitario es el proceso de evaluar la gravedad de la condición de un paciente y determinar el orden de atención médica en situaciones en las que los recursos son limitados. La implementación de modelos de aprendizaje automático en los centros hospitalarios permitiría ayudar en la clasificación y priorización de pacientes basándose en datos clínicos y síntomas relacionados con la enfermedad. Además, esto permitiría conocer de antemano la necesidad de recursos hospitalarios, como camas de cuidados intensivos o respiradores, y planificar en consecuencia.

El aprendizaje automático puede ser una herramienta útil en el triaje sanitario, pero no debe reemplazar la experiencia y el juicio clínico de los profesionales de la salud. Los modelos de aprendizaje automático deben utilizarse como una herramienta de apoyo a la toma de decisiones, pero siempre debe haber una supervisión y revisión médica adecuada para tomar decisiones finales.

El desarrollo de modelos de aprendizaje automático puede ser de gran utilidad en el triaje de pacientes y en la gestión de recursos en los centros hospitalarios durante la pandemia de la enfermedad. Estos modelos podrían ayudar a controlar la propagación de la enfermedad y mejorar la atención médica de los pacientes.

\subsection{Entendimiento de los datos}
La cohorte de datos han sido proporcionados por el Instituto de Investigación Sanitaria de Valdecilla (IDIVAL) sobre una población de 305 pacientes, cuyas muestras fueron recogidas entre Abril de 2020 y Marzo de 2021 y que tuvieron un resultado positivo en COVID-19. El valor de clasificación es el \textbf{score} de evolución clínica que tiene un valor de 1 a 5, siendo 1=asintomático; 2=tratamiento con gafas nasales; 3=tratamiento con ventimask; 4=Ingreso en UCI; 5=exitus.

El dataset se compone de un archivo excel con las siguientes características:
\begin{itemize}
	\item 420 filas o determinaciones.
	\item 97 columnas o marcadores.
    \item Sólo hay 158 filas con el valor de clasificación score cumplimentado.
\end{itemize}

Investigando los datos, vemos características que no vamos a utilizar debido a qué son fechas o números identificativos de muestras o pacientes que no nos aportan ningún tipo de información.

\subsection{Preparación de los datos}
Creamos un nuevo dataset copia del original y eliminamos las columnas que comentamos en la sección anterior que no nos aportan datos para clasificación.

Asimismo, observamos que el valor de clasificación \textit{score} tiene muchas filas vacías, lo que no nos permitiría realizar la clasificación, por lo que también las eliminamos.

Una vez tenemos las características con las que vamos a trabajar, observamos que tienen muchos valores perdidos, por lo que utilizaremos la \hyperref[imputadatos]{imputación de datos} para otorgar un valor aproximado para ese dato perdido.

\subsection{Modelado}
Aquí se aplican técnicas de modelado y análisis de datos para responder a las preguntas planteadas en la etapa de entendimiento del negocio. En el caso de este proyecto, se han realizado varios modelos, utilizando \hyperref[autosklearn]{Autosklearn} para realizar una búsqueda del clasificador que mejor funciona con el dataset. Los modelos generados con \hyperref[lda]{lda} han sido los dos métodos de clasificación con los mejores resultados.

\subsection{Evaluación}
En esta fase, se evalúan los modelos desarrollados. Se analizan los resultados y se comparan con el trabajo de partida inicial. Si los modelos no cumplen con los requisitos, se pueden realizar ajustes o mejoras.

\subsection{Despliegue}
Los modelos desarrollados se podrán poner a disposición de los profesionales para poder obtener un triaje de pacientes con síntomas COVID-19 a partir de una serie de muestras que se les recoge en el hospital.

\section{Imputación de datos} \label{imputadatos}
La imputación de datos se refiere al proceso de reemplazar los valores perdidos en un conjunto de datos con valores estimados o inferidos. Es una técnica comúnmente utilizada en el análisis de datos para abordar el problema de los valores perdidos, ya que estos pueden afectar la calidad y la validez de los resultados obtenidos.

Existen varios métodos y técnicas, pero en los modelos de este proyecto he utilizado dos:

\begin{itemize}
	\item \textbf{Imputación simple} Consiste en reemplazar los valores perdidos por un único valor, como la media, la mediana o el valor más frecuente del conjunto de datos. Esta técnica es rápida y sencilla, pero no tiene en cuenta la relación entre las variables y puede introducir sesgos en los resultados.
	\item \textbf{Imputación por vecinos más cercanos} Se basa en encontrar observaciones similares a aquellas con valores que faltan y utilizar los valores de las observaciones vecinas para estimar los valores perdidos. Esta técnica es útil cuando los datos tienen una estructura de vecindad o proximidad.
\end{itemize}
\capitulo{4}{Técnicas y herramientas}

Esta parte de la memoria tiene como objetivo presentar las técnicas metodológicas y las herramientas de desarrollo que se han utilizado para llevar a cabo el proyecto. Si se han estudiado diferentes alternativas de metodologías, herramientas, bibliotecas se puede hacer un resumen de los aspectos más destacados de cada alternativa, incluyendo comparativas entre las distintas opciones y una justificación de las elecciones realizadas. 
No se pretende que este apartado se convierta en un capítulo de un libro dedicado a cada una de las alternativas, sino comentar los aspectos más destacados de cada opción, con un repaso somero a los fundamentos esenciales y referencias bibliográficas para que el lector pueda ampliar su conocimiento sobre el tema.



\capitulo{5}{Aspectos relevantes del desarrollo del proyecto}

Este apartado pretende recoger los aspectos más interesantes del desarrollo del proyecto, comentados por los autores del mismo.
Debe incluir desde la exposición del ciclo de vida utilizado, hasta los detalles de mayor relevancia de las fases de análisis, diseño e implementación.
Se busca que no sea una mera operación de copiar y pegar diagramas y extractos del código fuente, sino que realmente se justifiquen los caminos de solución que se han tomado, especialmente aquellos que no sean triviales.
Puede ser el lugar más adecuado para documentar los aspectos más interesantes del diseño y de la implementación, con un mayor hincapié en aspectos tales como el tipo de arquitectura elegido, los índices de las tablas de la base de datos, normalización y desnormalización, distribución en ficheros3, reglas de negocio dentro de las bases de datos (EDVHV GH GDWRV DFWLYDV), aspectos de desarrollo relacionados con el WWW...
Este apartado, debe convertirse en el resumen de la experiencia práctica del proyecto, y por sí mismo justifica que la memoria se convierta en un documento útil, fuente de referencia para los autores, los tutores y futuros alumnos.

\capitulo{6}{Trabajos relacionados} \label{trabajosrel}

\section{Búsqueda de trabajos relacionados}
Investigación en los principales buscadores de artículos científicos acerca de documentos publicados relacionados con palabras como COVID-19, clasificación, aprendizaje automático, triaje y score. De entre todos, se discriminaron por los artículos que se basaran en \textit{datasets} que contuviesen marcadores recogidos a los pacientes para el diagnóstico del COVID-19.

En la tabla \ref{tabla:trabajosrelacionados}, se enumeran el listado de los trabajos relacionados con la fecha de publicación y una breve descripción.

\tablaConColores{Trabajos relacionados con este proyecto}{p{4cm} p{1.7cm} p{6cm}}{3}{trabajosrelacionados}
{Título & Fecha & Descripción\\}{ 
Innate and Adaptive Immune Assessment at Admission to Predict Clinical Outcome in COVID-19 Patients \cite{sansegundo:2021} & 29/07/2021 & Predicción de la severidad de la enfermedad en la admisión del paciente categorizada por los requisitos de oxigenoterapia\\
Machine-learning-based COVID-19 mortality prediction model and identification of patients at low and high risk of dying \cite{banoei:2021} & 08/09/2021 & Modelo de predicción de mortalidad para pacientes COVID-19 hospitalizados, así como una clasificación de pacientes para verificar los grupos de bajo y alto riesgo\\
Deep forest model for diagnosing COVID-19 from routine blood tests \cite{aljame:2021} & 17/08/2021 & Uso de técnicas de machine learning basadas en datos clínicos y/o de laboratorio para la detección temprana de COVID-19\\
Development of a severity of disease score and classification model by machine learning for hospitalized COVID-19 patients \cite{marcos:2021} & 21/04/2021 & Modelo de machine learning para identificar de manera temprana a los pacientes que morirán o requerirán ventilación mecánica durante la hospitalización a partir de las características clínicas y de laboratorio obtenidas en el momento del ingreso\\
Machine learning based predictors for COVID-19 disease severity \cite{patel:2021} & 25/02/2021 & Clasificación mediante random forest de la severidad en los pacientes con COVID-19 para calcular si necesitarán ventilación mecánica o/y ingreso en UCI\\
Prediction model and risk scores of ICU admission and mortality in COVID-19 \cite{zhao:2020} & 30/07/2020 & Desarrolla scores basados en las características clínicas en el momento de la admisión para predecir el ingreso y la mortalidad en UCI en pacientes con COVID-19\\
Scoring systems for predicting mortality for severe patients with COVID-19 \cite{shang:2020} & 15/07/2020 & Definición de sistemas de score para predecir el ingreso en UCI y/o la mortalidad en pacientes de COVID-19\\
A Multimodality Machine Learning Approach to Differentiate Severe and Nonsevere COVID-19: Model Development and Validation \cite{chen:2021} & 07/04/2021 & Uso de machine learning para diferenciar con precisión los tipos clínicos de COVID-19 graves y no graves en función de múltiples características médicas y proporcionar predicciones fiables del tipo clínico de la enfermedad\\
Multivariable mortality risk prediction using machine learning for COVID-19 patients at admission (AICOVID) \cite{kar:2021} & 17/06/2021 & Desarrollar y validar puntuaciones de riesgo de mortalidad individualizadas basadas en los datos clínicos y de laboratorio anonimizados en el momento del ingreso y determinar la probabilidad de Muertes a los 7 y 28 días\\
A systematic review of prediction models to diagnose COVID-19 in adults admitted to healthcare centers \cite{locquet:2021} & 18/06/2021 & Identificar, comparar y evaluar el desempeño de los modelos de predicción para el diagnóstico de COVID-19 en pacientes adultos en un entorno de atención médica\\
Machine learning is the key to diagnose COVID-19: a proof-of-concept study \cite{gangloff:2021} & 30/03/2021 & Desarrollar y evaluar modelos de aprendizaje automático utilizando datos clínicos y de laboratorio de rutina para mejorar el rendimiento de RT-PCR y CT de tórax para el diagnóstico de COVID-19 entre pacientes hospitalizados después de una urgencia\\
Comparing machine learning algorithms for predicting ICU admission and mortality in COVID-19 \cite{subudhi:2021} & 21/05/2021 & Compara el rendimiento de 18 algoritmos de aprendizaje automático para predecir la admisión y la mortalidad en la UCI entre pacientes con COVID-19\\
Rapid identification of SARS-CoV-2-infected patients at the emergency department using routine testing \cite{kurstjens:2021} & 29/06/2020 & Desarrollar un algoritmo para evaluar rápidamente el riesgo de una persona de infección por SARS-CoV-2 en el servicio de urgencias\\
Machine Learning for Mortality Analysis in Patients with COVID-19 \cite{sanchez:2021} & 12/11/2020 & Modelos de ML para encontrar reglas de decisión interpretables para estimar el riesgo de mortalidad de los pacientes que se pueden obtener del árbol de decisiones y que puede ser crucial en la priorización de la atención y los recursos médicos\\
}

\section{Clasificación de variables, técnicas y herramientas}

Una vez escogidos los trabajos sobre los que se basa la investigación para este proyecto, aparece la necesidad de clasificar de cada trabajo qué variables, técnicas y herramientas utilizaron los diversos autores en su investigación. De esta manera, se da la posibilidad de ver qué variables son las más usadas en las investigaciones, así como qué técnicas más se repiten y cuáles son las herramientas que se utilizan.

En la siguiente tabla \ref{tabla:clasificavariables} se expone qué variables con qué técnicas y herramientas se utilizaron en los trabajos relacionados.

\tablaConColores{Clasificación de variables, técnicas y herramientas}{p{5cm} p{3.4cm} p{3.4cm}}{3}{clasificavariables}
{Variables & Técnicas & Herramientas\\}{ 
Age, Ferritin, D-dimer, C4, \% de $CD8^{+}CD27^{-}CD28^{-}$, Suma absoluta de Linfocitos, \% de Monocitos no clásicos & Regresión Logística & GraphPad (Prism)\\
AST, WBC, Lymphocytes, Neutrophils, GGT, AGE, Basophils, Eosinophils, ALT, Platelets, gender, CRP, ALP, LDH, Monocytes & Extra trees, XGBoost, LightGBM, SHapley Additive exPlanations (SHAP) & Python (sklearn) \\
Old age, coronary heart disease (CHD), percentage of lymphocytes (LYM\%), procalcitonin (PCT), D-dimer (DD) & LASSO regression, Roc curve & Python\\
demographic variables (including age and sex), individual comorbidities and Charlson Comorbidity Index, chronic medical treatment, clinical characteristics, physical examination parameters, biochemical parameters & Random forest, Xgboost, Logistic regression, AUC & Python (sklearn, xgboost, eli5) \\
Leukocytes, Neutrophils, Lymphocytes, Eosinophil, Hemoglobin, Hematocrit, Platelets, ESR, BUN, Creatinine, Na, K, Ferritin, CRP, PCT, Lactate, Troponin, CK, BNP, LDH, Fibrinogen, ALT, AST, Albumin, D-dimer, Bilirubin, Prothrombin, APTT, pH, PaCo2, FiO2\_lab, Bicarbonate, CAD, Diabetes, Age > 65, AMS, Dementia, Nursing home, Q2 saturation < 88, yno2, Consolidation, Hypertension, Atrial fibrillation, Alcohol, Chest pain, Peripheral vascular disease, Stroke, Headache, Dyspnea, CRP, Lactate, Smoking & SIMPLS, AUC (Area under ROC curve), Decission Tree &	Python \\
age, sex, travel, contact history, and co-morbidities, fever, dyspnea, respiratory rate, and blood oxygen saturation (SpO2), RT-PCR, InterLeukin-6, D-Dimer, complete blood count, lipase, and C-reactive protein (CRP) & Random Forest, multilayer perceptron, support vector machines, gradient boosting, extra tree classifier, adaboost, Regresión, Validación cruzada quíntuple, AUC & Python \\
lactate dehydrogenase, procalcitonin, pulse oxygen saturation, smoking history, lymphocyte count, heart failure, chronic obstructive pulmonary disease, heart rate, age & Regresión Logística, AUC & TRIPOD \\
age, hypertension, cardiovascular disease, gender, diabetes, dimerized plasmin fragment D, high sensitivity troponin I, absolute neutrophil count, interleukin 6, lactate dehydrogenase & Random Forest, Predictive mean matching (PMM), Regresión Logística, AUC & R, Python \\
Age, Male Gender, Respiratory Distress, Diabetes Mellitus, Chronic Kidney Disease, Coronary Artery Disease, respiratory rate > 24/min, oxygen saturation below 90\%, Lymphocyte\% in DLC, INR, LDH, Ferritin & eXtreme Gradient Boosting (XGB), Random Forest, Regresión Logística, ‘Time to event’ using Cox Proportional Hazard Model, AUC ROC, Kaplan Meier (KM) Plots & Python \\
 & Prediction model study Risk of Bias Assessment Tool (PROBAST), XGBoost, Regresión Logísitica, AUROC & Python \\
Cough, Hyperthermia, Myalgias, Asthenia, Diarrhea, Confusion, Furosemid (usual treatment), Base excess, Lactates, Red blood cell count, Mean platelet volume, Leukocytes, Neutrophils, Platelet count, Eosinophils percentage, Basophils percentage, Lymphocytes, Monocytes, Ionogram, Potassium, Phosphor, Alanine aminotransferase, International normalized ratio, D-Dimer & binary logistic regression, random forest, artificial neural network, AUC, K-fold cross-validation & R-Studio (Dplyr, Purrr, missForest, Caret, corrplot, randomForest, neuralnet, pROC) \\
C-reactive protein, neutrophil percentages, lactate dehydrogenase, first respiratory, lower xygen saturation, lymphocyte percentages, estimated glomerular filtration rate (eGFR) < 60, high neutrophil percentage, high serum potassium, low lymphocyte percentages, high procalcitonin, D-dimer & AdaBoostClassifier, BaggingClassifier, GradientBoostingClassifier, RandomForestClassifier, XGBClassifier, ExtraTreesClassifier, LogisticRegression, DecisionTreeClassifier, LinearDiscriminantAnalysis, QuadraticDiscriminantAnalysis, MLPClassifier, PassiveAggressiveClassifier, Perceptron, LinearSVC, KNeighborsClassifier, GaussianNB & Python \\
Age, Gender, C-reactive protein, lactate dehydrogenase, ferritin, absolute neutrophil, lymphocyte counts & AUROC & Python, Excel 2010 \\
Patient ID, Age and gender, COVID diagnostic (confirmed/pending confirmation), ER date in, ER specialty, ER diagnostic, and destination after ER, First and last constant measurements in the ER (heart rate, temperature, minimum and maximum arterial pressure, O2 saturation in blood), Admission date to the hospital, ICU date in, ICU date out, and number of days in the ICU (if applicable), Discharge date and destination (home/deceased/transferred to other hospital/voluntary discharge/transferred to a socio-sanitary center) & Logistic regression, Survival Analisys, Decision tree, Random Forest, Bayessian networks, Biclustering, AUC & R (Caret, e1701, ggplot2), Python (sklearn, pandas, numpy  matplotlib) \\
}
\capitulo{7}{Conclusiones y Líneas de trabajo futuras}

Se exponen a continuación las conclusiones que se extraen de la realización de este proyecto, así como las posibles líneas de trabajo que se pueden considerar en un futuro.

\section{Conclusiones}

\subsection{Relacionadas con los resultados}

Los resultados obtenidos en los diferentes modelos realizados son parecidos a los obtenidos por el trabajo que se replicó. Esto se debe en parte a que se realizó sobre una pequeña cantidad de muestras y, por otra parte, a qué dichos datos no estaban del todo completos.

Estos resultados de los modelos se pueden mejorar con una mayor cantidad de muestras recogidas y con una mayor calidad en la recogida de los datos.

\subsection{Técnicas}
En este proyecto se ha realizado un seguimiento experimental riguroso:
\begin{itemize}
    \item Se realizó un estudio pormenorizado de los trabajos relacionados y la utilización que en ellos se hace de características, técnicas y herramientas, que deparó un mejor conocimiento de la realización de experimentos de aprendizaje automático con marcadores clínicos.
    \item Se prepararon los datos de la cohorte proporcionada, ya que contenían variables que no aportaban una información clara para nuestra investigación (como las fechas) y algunos valores nulos que fueron subsanados con la imputación de datos.
    \item La realización de la réplica del trabajo \textit{Innate and Adaptive Immune Assessment at Admission to Predict Clinical Outcome in COVID-19 Patients}\cite{sansegundo:2021} permitió conprender el porqué del trabajo original y el modelo de regresión logística aplicado.
    \item La utilización de autosklearn conllevó un estudio de esta biblioteca de python para poder entender qué valores introducir en el modelo automl que produjeran mejores resultados.
    \item Con la optimización de hiperparámetros con la biblioteca Optuna se pudo comprender cómo mejorar los experimentos en base a la realización de muchos otros experimentos variando los hiperparámetros.
    \item Conocimiento de biblioteca Shap para explorar la importancia de características del modelo de aprendizaje automático
\end{itemize}

\section{Líneas de trabajo futuras}

Como cualquier otro proyecto de índole experimental, una de las líneas de trabajo futuro es continuar con la experimentación. La tecnología avanza de forma imparable cada día y es por esto que, las bibliotecas científicas de software se actualizan y mejoran, las técnicas de aprendizaje automático se depuran y los sistemas dónde ejecutar las herramientas para los experimento no dejan de crecer en rendimiento. Todos estos factores hacen que un experimento realizado, pueda mejorarse de forma sustancial en el tiempo.

Otra línea clara a seguir es la distribución de los modelos para la utilización por parte de cualquier usuario, por ejemplo en un entorno web como se hizo en el artículo \textit{Development of a severity of disease score and classification model by machine learning for hospitalized COVID-19 patients} \cite{marcos:2021} con la distribución de la página \href{https://covid19salamanca-score.herokuapp.com/}{Herokuapp}; en esta aplicación, un profesional sanitario puede obtener una ayuda para identificar de manera temprana a los pacientes que morirán o requerirán ventilación mecánica durante la hospitalización a partir de las características clínicas y de laboratorio que debe introducir.


\bibliographystyle{plain}
\bibliography{bibliografia}

\end{document}
